% ---
% filename : ["electrotechnique_electronique_20260429_548_stabilisation_zener_charge.tex"]
% id : "electrotechnique_electronique_20260429_548"
% title : "Stabilisation de tension par diode Zener"
% domain : "Électrotechnique"
% subdomain : "Électronique"
% tags : ["diode zener", "stabilisation", "résistance", "puissance", "loi d'Ohm", "ts"]
% difficulty : 2
% time_solve : 20  # Calcul : 2 (difficulte) * (5 questions * 2)
% has_octave : true
% octave_files : ["electrotechnique_electronique_20260429_548_stabilisation_zener_charge.m"]
% author : "om"
% ---
\documentclass[a4paper,11pt,answers,addpoints]{exam}
% ---------------------------------------------------------
% LANGUE, ENCODAGE, FONTS
% ---------------------------------------------------------
\usepackage[utf8]{inputenc}
% ---------------------------------------------------------
% GÉOMÉTRIE ET MISE EN PAGE
% ---------------------------------------------------------
\usepackage[left=1.7cm,right=1.5cm,top=2cm,bottom=1cm]{geometry}
%\usepackage{a4wide}
\usepackage{microtype}      % Meilleure répartition du texte
\usepackage{float}          % Positionnement [H]
\usepackage{needspace}      % Saut conditionnel intelligent

% Éviter les veuves/orphelines (optionnel, à décommenter si besoin)
% \widowpenalty=10000
% \clubpenalty=10000

% Espacement vertical flexible
\raggedbottom
\setlength{\parskip}{0.5em plus 0.2em minus 0.1em}
\setlength{\parindent}{0pt}


% ---------------------------------------------------------
% MATHÉMATIQUES
% ---------------------------------------------------------
\usepackage{amsmath, amssymb, bm, esvect}

% Vecteurs
\newcommand{\V}{\overrightarrow}
\def\vect#1{\overrightarrow{\strut#1}}
\providecommand{\Degrees}[0]{\ensuremath{^\circ}}

% ---------------------------------------------------------
% UNITÉS ET GRANDEURS (SIunitx)
% ---------------------------------------------------------
\usepackage{siunitx}
\sisetup{output-decimal-marker={,}}
\DeclareSIUnit \var {var}
\DeclareSIUnit \VA {VA}
\DeclareSIUnit \kVA {kVA}
\DeclareSIUnit[quantity-product={}] \degree{\text{\textdegree}}

% ---------------------------------------------------------
% GRAPHIQUES : TikZ + CircuitikZ + PGFPlots
% ---------------------------------------------------------
\usepackage{tikz}
\usetikzlibrary{
	arrows.meta,
	intersections,
	decorations.pathreplacing,
	positioning
}

\usepackage[european,straightvoltages,EFvoltages]{circuitikz}

\usepackage{tkz-fct}
\usepackage{pgfplots}
\pgfplotsset{compat=1.12}

% Symbole étoile-triangle personnalisé
\newcommand{\wye}{%
	\mathbin{\tikz[x=1ex,y=1ex]{%
			\draw[line width=.1ex] (0,0)--(45:1)--++(-45:1) (45:1)--++(0,1);
	}}%
}


\usepackage{sankey}

\pgfdeclarelayer{background}
\pgfdeclarelayer{foreground}
\pgfdeclarelayer{sankeydebug}
\pgfsetlayers{background,main,foreground,sankeydebug}


% ---------------------------------------------------------
% TABLEAUX
% ---------------------------------------------------------
\usepackage{tabularx}
\usepackage{tabu}

% ---------------------------------------------------------
% HYPERLIENS
% ---------------------------------------------------------
\usepackage[colorlinks=true,
menucolor=red,
breaklinks=true,
urlcolor=blue,
linkcolor=blue]{hyperref}

% ---------------------------------------------------------
% COMMANDES PERSONNELLES
% ---------------------------------------------------------
\newcommand\nomauteur{bdminasmorgul@protonmail.com}
\newcommand\dateTe{29.04.2026}
\newcommand\brancheTe{TS}

% ---------------------------------------------------------
% STYLE DE L'EXERCICE
% ---------------------------------------------------------
\pagestyle{headandfoot}
\firstpageheadrule
\runningheadrule

\firstpageheader{\brancheTe\ (\nomauteur)}{Élève : }{(\dateTe) \thepage/\numpages}
\runningheader{\brancheTe\ (\nomauteur)}{Élève : }{(\dateTe) \thepage/\numpages}

% Compteur en bas de page
\newcommand{\totalpointtts}{%
	\begin{tikzpicture}[remember picture, overlay]
		\node[anchor=south west, inner sep=0pt, outer sep=0pt]
		at ([xshift=0.5cm,yshift=0.5cm]current page.south west)
		{\rotatebox{90}{Total des points : \numpoints}};
	\end{tikzpicture}
}

\firstpagefooter{\totalpointtts}{}{}
\runningfooter{\totalpointtts}{}{}

% Réglages exam
\printanswers
\addpoints
\pointsinmargin
\bracketedpoints

% ---------------------------------------------------------
% LISTINGS (code)
% ---------------------------------------------------------
\usepackage{xcolor}
\usepackage{listings}

\lstdefinestyle{octavestyle}{
	language=Matlab,
	basicstyle=\ttfamily\small,
	keywordstyle=\color{blue},
	commentstyle=\color{green!50!black},
	stringstyle=\color{red},
	stepnumber=1,
	frame=single,
	breaklines=true,
	breakatwhitespace=true,
	tabsize=4
}


% ---------------------------------------------------------
% DOCUMENT
% ---------------------------------------------------------
\begin{document}
\unframedsolutions
\pointsinmargin
\bracketedpoints

\section*{Préparation TS PQ CFC ELMO, IELE, PELE}

\begin{questions}
\question
On considère le montage ci-dessous (stabilisateur de tension) constitué d'une résistance série\\ $R_V = \qty{100}{[\Omega]}$, d'une diode Zener de tension $U_Z = \qty{6,2}{[V]}$ et d'une résistance de charge $R_L = \qty{220}{[\Omega]}$. La tension d'entrée est de $U_1 = \qty{12}{[V]}$. 

\begin{center}

	\begin{circuitikz}
		\draw (0,2) to[short, o-] (0.5,2);
		\draw (0,2) to[R, l^=$R_V$, current/distance = -1.7, i=$I_{Rv}$, o-] (3,2);
		\draw (3,2) to[short, -o] (5,2);
		\draw (3,0) to[zD*, l=$U_Z$, i^<=$I_Z$] (3,2);
		\draw (4.5,2) to[R, l_=$R_L$] (4.5,0);
		\draw (0,0) to[short, o-o] (5,0);
		\draw[->, >=stealth] (0,1.9) -- (0,0.1) node[midway, left] {$U_1$};
		\draw[->, >=stealth] (5,1.9) -- (5,0.1) node[midway, right] {$U_2$};
	\end{circuitikz}
\end{center}

\begin{parts}
	\part[3] Calculer la tension de sortie $U_2$ aux bornes de la charge $R_L$.
	\part[3] Déterminer l'intensité du courant $I_{Rv}$ circulant dans la résistance de protection $R_V$. 
	\part[3] Déterminer l'intensité du courant $I_L$ absorbée par la charge. 
	\part[3] En déduire l'intensité du courant $I_Z$ qui traverse la diode Zener. 
	\part[3] Calculer la puissance $P_Z$ dissipée par la diode Zener. 
\end{parts}
\vspace{0.5em}
\needspace{55\baselineskip}
\begin{solution}
	\begin{align*}
		a) \quad \text{Comme } &U_1 > U_Z \text{, la diode Zener est passante et stabilise la tension.}\\
		U_2 = U_Z = \qty{6,2}{[V]} \\[6pt]
		b) \quad I_{Rv} &= \dfrac{U_1 - U_2}{R_V} = \dfrac{12 - 6,2}{100} = \qty{0,058}{[A]} = \qty{58}{[mA]} \\[6pt]
		c) \quad I_L &= \dfrac{U_2}{R_L} = \dfrac{6,2}{220} = \qty{0,0282}{[A]} = \qty{28,2}{[mA]} \\[6pt]
		d) \quad I_Z &= I_{Rv} - I_L = 58 - 28,2 = \qty{29,8}{[mA]} \\[6pt]
		e) \quad P_Z &= U_Z \cdot I_Z = 6,2 \cdot 0,0298 = \qty{0,1848}{[W]} = \qty{184,8}{[mW]} \\
	\end{align*}

	\begin{verbatim}
		% Télécharger OCTAVE depuis https://wiki.octave.org/Using_Octave et l'installer
		% Puis copier-coller le code dans OCTAVE
		
		U1 = 12;    % [V] Tension d'entrée
		Rv = 100;   % [Ohm] Résistance série
		Uz = 6.2;   % [V] Tension Zener
		Rl = 220;   % [Ohm] Résistance de charge
		
		% a) Tension de sortie
		U2 = Uz;
		printf("a) Tension de sortie : U2 = %.2f [V]\n", U2);
		
		% b) Courant dans la résistance série
		Irv = (U1 - U2) / Rv;
		printf("b) Courant Irv : Irv = %.4f [A] (%.2f [mA])\n", Irv, Irv*1000);
		
		% c) Courant dans la charge
		Il = U2 / Rl;
		printf("c) Courant charge Il : Il = %.4f [A] (%.2f [mA])\n", Il, Il*1000);
		
		% d) Courant dans la Zener
		Iz = Irv - Il;
		printf("d) Courant Zener Iz : Iz = %.4f [A] (%.2f [mA])\n", Iz, Iz*1000);
		
		% e) Puissance Zener
		Pz = Uz * Iz;
		printf("e) Puissance dissipee Pz : Pz = %.4f [W] (%.2f [mW])\n", Pz, Pz*1000);
	\end{verbatim}
	\end{solution}
\end{questions}
\end{document}